\documentclass{article}
\usepackage{geometry}
\geometry{
	a4paper,
	left=2.5cm,      % 左边距 2cm
	right=2.5cm,     % 右边距 2cm
	top=2.5cm,       % 上边距 2cm
	bottom=2.5cm,    % 下边距 2cm
	headheight=13.6pt
}
\usepackage{amsmath}
\usepackage{booktabs} % 用于三线表
\usepackage{caption}
\usepackage{setspace} % 用于设置行间距
\usepackage{array}
\usepackage{CTeX}
\setlength{\headheight}{13.6pt}
\usepackage{enumitem} % 用于自定义列表格式

\begin{document}
	\subsection{利用统计量描述xxx}
	
	\subsubsection{描述统计量}
	
	本文引入统计学和数据分析中的三个描述性统计量：\textbf{变异系数}、\textbf{偏度系数}和\textbf{峰度系数}分析数据的分布规律。变异系数用于衡量数据的\textbf{离散程度}，偏度系数用于衡量数据的\textbf{分布形态的对称性}，峰度系数用于衡量数据的\textbf{分布形态的陡峭程度} \cite{1}。以下为三个系数的具体介绍：
	
	\begin{enumerate}
		\item \textbf{变异系数}：
		
		变异系数是一个\textbf{无量纲}的统计量，用于比较不同数据集的离散程度，还可以作为不同量纲之间的桥梁，例如它可以消除身高和体重平均值，标准差等的差异，进行有意义的比较，通常使用百分比形式表示，计算公式如下：
		
		\begin{equation}
			CV = \frac{\sigma}{\mu}\times100\%
			\label{eq:cv}
		\end{equation}
		
		其中，\(CV\)是变异系数，\(\sigma\)是总体标准差，\(\mu\)是总体均值。CV值越大，表示数据的\textbf{相对波动性越大}；CV值越小，则相对波动性越小，\textbf{数据越集中} \cite{2}。
		
		\item \textbf{偏度系数}：
		
		\textbf{偏度系数}是衡量数据分布\textbf{不对称性}程度的统计量，描述分布偏离对称分布的程度和方向，计算公式如下：
		
		\begin{equation}
			SK = \frac{\sum_{i=1}^{m} (x_{ij} - \bar{x})^3 / m}{s_j^3}
			\label{eq:sk}
		\end{equation}
		
		其中，\(m\)是样本数量，\(x_{ij}\)是每个数据点，\(\bar{x}\)是样本均值，\(s_j\)是样本标准差，\(i\)是数据点，\(j\)是度量项目。若\textbf{SK = 0}，则表示数据分布是\textbf{对称}的；若\textbf{SK > 0}，则表示正偏态或右偏态，即数据点落在\textbf{右侧偏多}，均值 > 中位数 > 众数；若\textbf{SK < 0}，则表示负偏态或左偏态，即数据点落在\textbf{左侧偏多}，均值 < 中位数 < 众数，偏度系数的\textbf{绝对值}越大，表示\textbf{偏斜程度}越严重 \cite{3}。
		
		\item \textbf{峰度系数}：
		
		\textbf{峰度系数}是衡量数据分布\textbf{陡峭程度}的统计量，衡量数据分布与正态分布相比最高点的高度以及低点的高度，计算公式如下：
		
		\begin{equation}
			K = \frac{\sum_{i=1}^{m} (x_{ij} - \bar{x})^4 / m}{s_j^4} - 3
			\label{eq:k}
		\end{equation}
		
		其中，\(K\)是峰度系数，\(m\)是样本数量，\(x_{ij}\)是每个数据点，\(\bar{x}\)是样本均值，\(s_j\)是样本标准差，\(i\)是数据点，\(j\)是度量项目。若\textbf{峰度 = 0}，则表示分布与正态分布的陡峭程度相同，为\textbf{常峰态}；若\textbf{峰度 > 0}，则表示数据分布比正态分布\textbf{更陡峭}，\textbf{尾部更厚}，数据中有更多的\textbf{极端值}远离均值；若\textbf{峰度 < 0}，则表示数据分布比正态分布更\textbf{平坦}，\textbf{尾部更薄}，数据更加\textbf{分散}，极端值较少 \cite{3}。
	\end{enumerate}
	
	\subsubsection{得出统计结果}
	
	本文利用MATLAB求解上述三种统计量，由于篇幅原因，仅展示部分单品结果，全部结果见文件Statistics.xlsx：
	
	\renewcommand{\arraystretch}{1.5}
	\begin{table}[h!]
		\centering
		\caption{表5 预测结果}
		\begin{tabular}{cccc}
			\toprule[1.5pt]
			\textbf{单品组合} & \textbf{2021年相关系数} & \textbf{2022年相关系数} & \textbf{2年平均相关系数} \\
			\midrule
			上海青螺丝椒(份) & -0.410018296 & -0.903074635 & -0.656546465 \\
			上海青小青菜(份) & -0.426960775 & -0.884997079 & -0.655978927 \\
			云南生菜螺丝椒(份) & -0.468748579 & -0.83800917 & -0.653378874 \\
			上海青云南生菜(份) & -0.429652604 & -0.871453397 & -0.650553 \\
			上海青杏鲍菇(2) & -0.404275203 & -0.894663846 & -0.649469525 \\
			上海青奶白菜(份) & -0.426960775 & -0.86276201 & -0.644861392 \\
			上海青菜心(份) & -0.426960775 & -0.845851263 & -0.636406019 \\
			\bottomrule[1.5pt]
		\end{tabular}
	\end{table}
	
	接下来以三种不同系数进行分类分析：
	
	\begin{enumerate}
		\item \textbf{变异系数}：由于变异系数多为负数，剩下的都大于0.15属于强变异，在统计学意义上无很大意义，故不做分析。
		\item \textbf{偏度系数}：风化铅钡类别的有9组成分指标的偏度系数小于0，说明落在均值左侧的数据均偏多，其中Na2O偏度系数近似等于0，说明CaO的数据是接近对称分布的，其余五组成分指标的偏度系数大于0，说明落在均值右侧的数据偏多。
		\item \textbf{峰度系数}：六组成分指标的峰度系数大于0，说明指标分布相比于正态分布顶部更加尖锐或者尾部更加粗，其余八项指标峰度系数小于0，说明分布相比于正态分布顶部更加平坦或者尾部更加细。风化高钾类别的偏度系数六组小于0，说明该六组落在均值左侧的数据均偏多，还有六组偏度系数值为0，说明该六组数据是对称分布的，其余两组Al2O3、P2O5偏度系数值大于0，数据分布在均值右侧的数据偏多；六组数据峰度系数大于0，说明指标分布相比于正态分布顶部更加尖锐或者尾部更加粗，还有两组数据峰度系数小于0，说明分布相比于正态分布顶部更加平坦或者尾部更加细，其余六组数据偏度系数为0，说明其分布与标准正态分布顶部尖锐程度和尾部粗细相当。
	\end{enumerate}
	
	
		%\bibitem{1} 苏为华.多指标综合评价理论与方法问题研究[D].厦门大学,2000.
		%\bibitem{2} 王春枝,斯琴.德尔菲法中的数据统计处理方法及其应用研究[J].内蒙古财经学院学报(综合版),2011,9(04):92-96.DOI:10.13895/j.cnki.jimufe.2011.04.019.
		%\bibitem{3} 李兴奇,高晓红.无量纲化方法的有效性评价[J].统计与决策,2021,37(15):24-28.DOI:10.13546/j.cnki.tjyjc.2021.15.005.
	
\end{document}